% ============================================================
% Conclusiones — cierre del Volumen I
% Traducción de "Théorie mathématique de la lumière" (1889), "CONCLUSIONS"
% Original: §§241-243, páginas impresas 398-401
% ============================================================
\chapter*{Conclusiones}
\addcontentsline{toc}{chapter}{Conclusiones}
\markboth{Conclusiones}{Conclusiones}

\noindent\textbf{241.} En el estudio de cada fenómeno hemos expuesto
paralelamente varias teorías que dan cuenta igualmente bien de los hechos
observados. Estas teorías pueden, por lo demás, agruparse en uno de los
dos grupos siguientes: aquellas en las que se supone, como Fresnel, que la
elasticidad $\varepsilon$ del medio es constante; aquellas en las que se
admite, con Neumann, que la densidad $\rho$ del éter es constante. No
hemos encontrado ninguna razón que pueda hacer preferir una de estas
hipótesis a la otra. Solo la explicación de la aberración por el arrastre
parcial del éter puede hacer inclinar la balanza del lado de la hipótesis
de Fresnel; pues el arrastre parcial supone que la densidad del éter no es
la misma en todos los medios. Pero, como lo hemos hecho notar, es difícil
darse bien cuenta de estos fenómenos de aberración y ninguna teoría es
satisfactoria. No hay, pues, ahí razón suficiente para decidir la elección
de una teoría.

Por lo demás, no podemos quejarnos de estar en la imposibilidad de hacer
una elección. Esta imposibilidad nos muestra que las teorías matemáticas de
los fenómenos físicos no deben ser consideradas más que como instrumentos
de investigación; instrumentos muy preciosos, es cierto, pero de los que no
debemos quedar esclavos y que debemos rechazar en cuanto se encuentren en
contradicción formal con la experiencia.

\noindent\textbf{242.} Hay una razón general que nos impide elegir entre
las teorías ópticas que hemos expuesto. Sabemos, en efecto, que las
ecuaciones del movimiento en un medio elástico isótropo o anisótropo son
ecuaciones lineales y de coeficientes constantes. Una propiedad general de
las ecuaciones de este género es que, si $\xi_1$ y $\xi_2$ son dos
integrales de una de ellas, la cantidad $A\xi_1+B\xi_2$ será también una
solución. Tenemos, pues, una infinidad de maneras de satisfacer a los
problemas ópticos.

Además hemos visto que una de las ecuaciones del movimiento en un medio
elástico isótropo es\footnote{La numeración de ecuaciones se reinicia una
vez más, ya dentro de las Conclusiones del volumen.}
\setcounter{equation}{0}
\begin{equation}
\frac{d^2\xi}{dt^2} = \Delta\xi - \frac{d\Theta}{dx}.
\end{equation}
Si derivamos los dos miembros de esta ecuación respecto de una variable
cualquiera, tendremos
\[
\frac{d^2\xi'}{dt^2} = \Delta\xi' - \frac{d\Theta'}{dx}.
\]
Luego, si una función satisface a la ecuación (1), una derivada cualquiera
de esta función satisfará también. Si designamos por $\xi, \eta, \zeta$
las componentes del desplazamiento de una molécula de éter en la teoría de
Fresnel, cantidades que satisfacen las ecuaciones del movimiento, tales
como (1), las derivadas $\xi_x', \xi_y', \xi_z', \eta_x', \ldots$
satisfarán también a las ecuaciones. Es fácil comprobar que ocurrirá lo
mismo con los binomios alternados $\zeta_y'-\eta_z'$, $\xi_z'-\zeta_x'$,
$\eta_x'-\xi_y'$.

Ahora bien, sabemos que las componentes del desplazamiento de una molécula
de éter en la teoría de Neumann son precisamente esos binomios alternados;
por consiguiente, serán también soluciones de las ecuaciones del
movimiento. Se puede, pues, estar seguro de que todo fenómeno explicado por
la teoría de Fresnel lo será igualmente por la de Neumann. La recíproca es,
por lo demás, cierta.

En algunos casos, sin embargo, hemos tenido coeficientes variables en las
ecuaciones del movimiento. Así en la teoría de M.~Sarrau y en la
consideración de la capa de paso en la reflexión. Pero, como ignoramos
absolutamente la ley de variación de estos coeficientes, nos basta con
algunas hipótesis para hacer concordar la teoría con la experiencia y no
podemos decidir nada sobre la exactitud de la teoría.

\noindent\textbf{243.} El estudio de un fenómeno particular, la
\emph{polarización por difracción}, parecía poder permitir decidir entre
la teoría de Fresnel y la de Neumann. Se creía haber llegado, por el
cálculo, al resultado de que la rotación del plano de polarización no
tenía el mismo valor en estas dos teorías. Pero ciertamente los cálculos
eran inexactos, pues, siendo las ecuaciones del movimiento de coeficientes
constantes, no puede, como acabamos de decir, haber ninguna divergencia
entre los resultados de las dos teorías. Sin embargo, a raíz de estos
cálculos, se intentaron experiencias en esta vía. Son muy delicadas, pues
la desviación de un rayo difractado es muy débil y la rotación del plano de
polarización que de ella resulta es excesivamente pequeña. Con el objeto de
aumentar la desviación del rayo difractado, se ha experimentado con redes.
Pero el problema se complica entonces, pues hay a la vez difracción y
refracción o reflexión; así, las experiencias no estuvieron de acuerdo
ni entre ellas, ni con las consecuencias de la teoría de Fresnel, ni con
las de la teoría de Neumann.

Muy recientemente el señor Gouy retomó el mismo problema sin idea
preconcebida; obtuvo una desviación considerable del rayo difractado
colocando la fuente luminosa, formada por la concentración de rayos en el
foco de una lente convergente, sobre el borde de una pantalla. Los
resultados de sus experiencias para el valor de la rotación del plano de
polarización tampoco están de acuerdo con las dos teorías de la
polarización. La teoría de la difracción se encuentra incluso en defecto,
pues el señor Gouy comprobó que los fenómenos dependen de la forma del
borde de esa pantalla y de la naturaleza de esta. Este desacuerdo entre las
experiencias del señor Gouy y la teoría de Fresnel no debe sorprendernos,
pues hemos dicho que era imposible encontrar una solución de la ecuación
\[
\Delta\xi + \alpha^2\xi = 0
\]
que satisficiera \emph{exactamente} a las condiciones del problema. Solo
satisfaciéndola \emph{aproximadamente} hemos podido edificar una teoría de
la difracción.

La aproximación era muy ampliamente suficiente en las condiciones
habituales de las experiencias de difracción; pues las cantidades
despreciadas son entonces extremadamente pequeñas. Ya no ocurre lo mismo en
las condiciones en que se situó el señor Gouy.

\bigskip
\begin{center}
\textsc{Fin}
\end{center}
